\section{本章小结}
\label{sec:smap-summary}

本章开头问：只知道相对距离，能否安排坐标；知道两组对应点，能否求出它们之间的映射。两个问题的第一步相同，都是辨认观测不能确定的东西。距离对整体平移、旋转和反射一无所知，因此最多把点配置确定到合同为止；缺失的距离还会留下真正的形变，例如只测边长的正方形可以压成菱形。固定参照能消去合同的自由度，却补不出缺失的观测；用仪器测得的锚点则是额外的观测，不是约定。

完整距离表的重构由双中心化完成。平方距离经过$-\frac12\bm J\bm\Delta\bm J$变成以质心为原点的Gram矩阵，它半正定当且仅当距离表欧氏可实现，秩就是所需的最小维数，谱分解直接给出坐标。完整距离相同的两组点必定合同，所以合同是完整表下唯一的不确定性。指定低维时，保留最大的正特征值在Gram矩阵的Frobenius误差下最优，负谱只会增加误差；但距离误差是另一个量，应力准则下截断坐标一般不是最优的。

部分距离的问题由观测图描述。连通不保证形状确定，局部刚性也不保证全局唯一，两个共边三角形就有互为反射的两个实现。缺边、不相容与含噪是三种不同的困难，残差只能建在已观测的边上；把缺失填成零等于断言两点重合，会让谱重构给出一个看似确定的错误形状。稳定性方面，完整表的误差经双中心化至多减半地传到Gram矩阵，特征值稳定，坐标轴方向取决于谱隙，小特征值经开方后最为脆弱；部分距离的稳定性则取决于约束本身，与完整表的谱隙无关。

映射估计先固定模型族，因为有限个点对确定不了任意的光滑映射。刚体模型用质心消去平移，再由交叉矩阵的奇异值分解求旋转；允许反射的正交对齐用于比较重构坐标，带行列式修正的旋转对齐才对应刚体位姿。点重合或共线时交叉矩阵秩不足，最优旋转成族出现，零残差并不说明变换已被确定。仿射模型用增广坐标化为线性最小二乘，源点仿射张成定义域时参数唯一，所得的$\bm A$是否可逆则要另行核验。

射影模型在齐次坐标下只确定到一个非零倍数。四对处于一般位置的点使直接线性变换的核恰为一维，且对应可逆矩阵；点共线时齐次解可能全是奇异的，必须核验。含噪时先作坐标归一化，再用单位范数的齐次最小二乘求代数解；代数残差是以$c_i^2$加权的重投影误差，依赖于坐标表示，代数最优不等于几何最优。几种模型共同说明：残差小、方程有解与参数唯一、结果稳定是不同的结论，要分别核对。

本章的点之间要么有距离，要么有对应，对象始终是有限个点本身。第~\ref{chap:manifold-recovery}章把有限个点看成一个未知流形上的样本，研究怎样由它们恢复切方向、局部曲面、沿流形的距离以及同调，那里的距离图与样本复形都只用到本章意义下的成对距离。
